\section{Checklist de preparação}

\subsection{Antes de qualquer protocolo}

\begin{longtable}{p{0.08\textwidth}p{0.84\textwidth}}
\textbf{OK} & \textbf{Item} \\
\hline
$\square$ & Confirmar com a coordenação do curso qual versão da Resolução 01/2003 deve ser aplicada ao caso. \\
$\square$ & Confirmar se o depósito de patente é exigido literalmente ou se registro de software/produto técnico será aceito. \\
$\square$ & Definir autores, inventores, titulares e eventual participação institucional. \\
$\square$ & Verificar regras internas de propriedade intelectual da instituição e comunicação ao NIT. \\
$\square$ & Fechar versão do SIGEV que será protegida. \\
$\square$ & Guardar instalador, código-fonte e documentação da versão fechada. \\
\end{longtable}

\subsection{Para registro de software}

\begin{longtable}{p{0.08\textwidth}p{0.84\textwidth}}
\textbf{OK} & \textbf{Item} \\
\hline
$\square$ & Gerar pacote compactado do código-fonte e documentação. \\
$\square$ & Calcular \textit{hash} SHA-256 do pacote. \\
$\square$ & Preencher dados de autor(es), titular(es), linguagem, campo de aplicação e data de criação. \\
$\square$ & Obter certificado digital qualificado ICP-Brasil. \\
$\square$ & Emitir e pagar GRU de registro de programa de computador. \\
$\square$ & Assinar Declaração de Veracidade. \\
$\square$ & Protocolar no e-Software. \\
$\square$ & Guardar certificado de registro e comprovantes. \\
\end{longtable}

\subsection{Para depósito de patente}

\begin{longtable}{p{0.08\textwidth}p{0.84\textwidth}}
\textbf{OK} & \textbf{Item} \\
\hline
$\square$ & Fazer busca de anterioridade em bases de patentes e artigos. \\
$\square$ & Revisar patenteabilidade com NIT/procurador. \\
$\square$ & Completar relatório descritivo com exemplos, telas e fluxos técnicos. \\
$\square$ & Revisar reivindicações para evitar proteção indevida de software em si ou método abstrato. \\
$\square$ & Preparar resumo em até formato aceito pelo INPI. \\
$\square$ & Avaliar necessidade de desenhos/fluxogramas. \\
$\square$ & Emitir e pagar GRU de patente. \\
$\square$ & Protocolar no e-Patentes. \\
$\square$ & Acompanhar exigências formais, publicação, pedido de exame e anuidades. \\
\end{longtable}

\subsection{Arquivos sugeridos no dossiê acadêmico}

\begin{itemize}
  \item \texttt{dist/SIGEV-Setup-x64.exe}: instalador final.
  \item \texttt{README.md} e \texttt{docs/}: documentação do aplicativo.
  \item \texttt{documentos\_inpi/pacote\_inpi\_sigev.tex}: pacote LaTeX revisável.
  \item PDF compilado do pacote LaTeX.
  \item Comprovante de protocolo ou certificado emitido pelo INPI.
  \item Declaração do orientador/coordenação indicando relação do software/patente com a dissertação.
\end{itemize}
